\section{Plan de capacitación y gestión del cambio}

\subsection{Objetivo de la gestión del cambio}

La estrategia de capacitación tiene como propósito facilitar la transición tecnológica de Machu Picchu Exclusive Tours E.I.R.L. desde una gestión manual basada en WhatsApp y Excel hacia un entorno digital centralizado en Odoo. Se busca reducir la brecha de aprendizaje de la gerencia y el personal operativo (incluyendo el apoyo familiar mencionado), asegurando que el sistema sea percibido como una herramienta que garantiza la seguridad de la información confidencial (pasaportes) y facilita el seguimiento de itinerarios complejos de larga duración.
\\
\\
Este enfoque coincide con lo establecido en el documento de Alcance, donde se señala la necesidad de capacitar a un equipo de usuarios clave en el flujo \emph{Oportunidad a Cierre}, en la gestión de la ficha del cliente y en el uso de actividades, como condición indispensable para que los módulos de CRM y Proyectos sean utilizados de manera correcta y sostenible en el tiempo.

\subsection{Metodología de transferencia de conocimientos}

El plan de capacitación se estructura en tres etapas fundamentales para asegurar la sostenibilidad del proyecto:

\begin{itemize}
    \item[a.] \textbf{Capacitación operativa (CRM)}: entrenamiento en el flujo integral de \emph{Oportunidad a Cierre}, incluyendo el registro sistemático de pasajeros, la carga segura de documentos de identidad en el repositorio estándar y la gestión de notas para mantener una visión 360° del turista extranjero.
    \item[b.] \textbf{Gestión de reservas e itinerarios (Proyectos)}: instrucción en el uso del tablero Kanban para el control de servicios (de 2 a 20 días), carga de vouchers de hoteles y boletos de tren, y el registro manual de adelantos para el control de saldos antes de compras no reembolsables.
    \item[c.] \textbf{Gestión de actividades y seguimiento}: taller práctico sobre el uso de alertas y recordatorios para llamadas y correos de fidelización, permitiendo una atención personalizada post-viaje.
\end{itemize}

Para reforzar el aprendizaje, se entregarán videos tutoriales operativos como material de consulta permanente.

\subsection{Despliegue y retroalimentación}

El cronograma de capacitación (3 semanas) se integra con las fases finales de configuración de Odoo. Esto permite realizar sesiones de retroalimentación con la gerencia para ajustar las etapas del embudo de ventas y las vistas de calendario antes de la puesta en marcha definitiva, validando que el sistema se adapte al flujo real de atención de clientes internacionales.

\subsection{Resultados esperados de la adopción tecnológica}

Dado que la etapa de capacitación se encuentra en curso o recién concluida, los resultados concretos de la adopción aún no pueden cuantificarse. Se espera que, una vez consolidado el uso del sistema por parte del personal, se observe una mejora gradual en la centralización de la información, el uso del CRM para el seguimiento de clientes y el control de reservas, así como una mayor adherencia a los protocolos de seguridad de la información. La evaluación de estos aspectos se realizará en la fase de monitoreo y ajustes posteriores.

